\section{En deterministisk felmodell}
\label{sec:errors}

En buss som alltid levererar varje byte i rätt ordning är inte en buss, utan ett funktionsanrop med
extra steg. Verkliga bussar tappar bytes, ändrar bitar och fördröjer leveranser.

Problemet med att testa felhantering är att verkliga fel är slumpmässiga. Ett test som ibland
fallerar säger ingenting, och ett fel som inträffade en gång går inte att återskapa.

\textbf{Lösningen är att göra felen deterministiska.} Bussimplementationen i kursens simulering
räknar bytes och gör exakt samma sak varje körning:

\begin{cppcode}
// Deterministic fault injection: byte five is dropped, every run, every time.
if (5U == ++myByteCounter) { return; }
\end{cppcode}

Tre fel räcker för att täcka det mesta:

\begin{booktable}{@{}lL@{}}
\thead{Fel} & \thead{Vad som händer} \\ \midrule
\textbf{Tappad byte} & Alla följande fält förskjuts. Parsern läser fel längd, väntar på fel antal
  bytes, och faller till slut på checksumman. \\
\textbf{Ändrad bit} & Framen är strukturellt korrekt --- rätt längd, rätt fältplacering --- men
  checksumman går inte ihop. Det är precis det fall checksumman finns för. \\
\textbf{Fördröjning} & Framen kommer fram, hel och riktig, men senare än väntat. Ingenting går
  sönder, och det är poängen: en korrekt parser bryr sig inte om \emph{när} bytesen kommer. \\
\end{booktable}

Det tredje fallet är värt att stanna vid. En parser som fungerar när bytesen kommer tätt men går
sönder när de kommer glest har ett dolt tidsberoende, och det är en bugg som aldrig visar sig
förrän systemet flyttas till långsammare hårdvara.
